% Discussion (Task 46, Scientific Reports): v3 Discussion, Limitations and Conclusion folded into one section without
% subheadings (plan 42, section 1.3). The limitation detail that does not stay here is Supplementary Note 5.
\section*{Discussion}
\label{sec:discussion}

On both datasets, rollout attention agreed weakly to moderately with attribution methods and kept most of its
top-\paramTopK{} mass under the feature constraints, yet it was less faithful than IG and occlusion by deletion and
insertion AUC.
An analyst who reads rollout attention as feature importance on these models would therefore see a stable picture that
only partly reflects what drives the prediction; when a decision rests on feature importance, IG or occlusion were the
more faithful choice on both datasets.
The stiffness of attention is a statement about the constraints and budgets we used, not a security property: among
the clean top-\paramTopK{} IG features that A2 targets, on average \topImmIGSeventeen{} on 2017 and \topImmIGEighteen{}
on 2018 were immutable, so the stiffness is partly structural, and the unconstrained reference moved attention
further.
The faithfulness loss after an evasion appeared for most explanation methods, including static class-mean rankings that
do not depend on the input, so any explanation of a row that an attack has flipped is less faithful than on clean
traffic; a plausible reading is that adversarial rows sit where the model's output depends on the features
differently, but this mechanism was not tested.

The consistency scores helped against an attacker that ignored them, but that gain was not specific to attention:
controls that used other attribution methods instead of attention, or statistics of the logits and distances, obtained
most or all of it, and the shuffled-attention control a1 still carries the class information of the true label.
The small gain that D$_+$ showed against A8 did not hold against A8x, which uses exact IG in the loop and selects
iterates by the true score: the sign reversed on 2017 in the pooled estimate and in all but
\axHeadOppSeedsSeventeen{} of the \nSeeds{} seeds, and giving D's attacker the same selection changed the outcome of
\sxOutcomeChangedPhrase.
An evaluation with stand-ins alone would therefore have overstated what D$_+$ achieves; this is the general point made
about adaptive attacks on detectors of adversarial examples \cite{tramer2020}, shown here for consistency monitoring.
We did not test why D$_+$ was easier to evade than D.
One untested reading is that each added score gives the attacker a further direction to move in while the threshold
stays at \paramFPR{} false-positive rate (FPR); the control a1, under its own stand-in attacker, also left more rows
undetected than D on 2017, but a1 uses the true class label and a weaker attacker, so it does not test this reading.
For a defender, we do not recommend these consistency scores as an addition to a confidence-and-distance detector
against adaptive attackers: they helped only against an attacker that ignored them, that help was available without
attention, and against an attacker that modelled them exactly they made evasion easier on both datasets, on 2018
close to the floor.
For anyone evaluating an explanation-based attack detector, the attacker should optimise the detector's exact score
wherever a gradient exists and select its iterates by the true score; an evaluation with differentiable stand-ins alone
can give a result of the opposite sign, as it did here.

The difference between the datasets in detector strength was larger than the difference between D and D$_+$, and on
each dataset a single class dominated the successful and undetected rows.
Conclusions about learned attack detectors on these benchmarks should therefore be reported per dataset and per class,
and on 2018 the detector comparisons have little room to differ, because D alone left only \axHeadVoneEighteenD{} of
the source rows successful and undetected.
Near duplicates between the training and test sets were common in both datasets and made the affected rows easier for
both detectors on 2017; removing them did not change any conclusion here, but evaluations on these benchmarks that keep
them may report detectors as stronger than they are on new traffic.
The headline difference under the exact-objective attacker kept its sign on all three row sets, whereas the difference
under the stand-in attacker changed more with the row set, which is a further reason to read the stand-in result as
fragile.

Several limitations bound these conclusions.
The results rest on \nSeeds{} seeds, and in the equal-effort comparison against A8x \axHeadOppSeedsPhrase{} had the
opposite sign; our bootstrap intervals resample source rows within seed, reuse the same source rows in every seed and
do not capture seed variance, so they do not show that an effect holds across seeds.
A8x used surrogate gradients, soft ranks with an annealed temperature, for the \nRankScores{} rank-based consistency
scores, which have no informative exact gradient; a better surrogate or a gradient-free search on these terms could
only raise the share for D$_+$.
The equal-effort protocol equalises candidates, steps and settings, not optimiser quality, and whether A6 is as strong
an attack on D as A8x is on D$_+$ was not measured; A6x matched the selection rule, not the optimiser.
The attackers of the controls a1 and a2 used stand-ins like A8, so their results may understate how easily these
controls can be evaded, and their comparison with D$_+$ under A8x is not at equal attacker strength.
The attribution-only control b was not attacked adaptively, because its attack run exceeded the memory of the graphics
card, so its family is untested against an adaptive attacker.
All attacks changed flow features, not packets, and a feature vector that our checker accepts may still not correspond
to traffic that an attacker can produce; the attackers knew the weights of the model and, where stated, of the
detector.
The attackers A8x and A6x were designed after the earlier results were known; their outcome rules and both amendments
were recorded before any of their attack rows existed.
We studied two datasets and one model family, without human evaluation or deployment, and both datasets carry known
design and labelling issues of benchmark NIDS data \cite{flood2024,lanvin2022}.
Further limitations on the constraint model, the detectors and the data are in Supplementary
Note~\suppNoteLimits.

In summary, attention in this flow classifier was stable under protocol-constrained attacks but only a partial
explanation, and attention-attribution consistency scores did not make a confidence-and-distance detector harder to
evade for attackers that optimised its exact score.
Evaluations of explanation-based attack detectors need such attackers; on our two datasets the consistency scores did
not pass that test.
The attribution-only control under an adaptive attacker, the rank-based scores under an exact search and the relative
optimiser strength of the attackers remain open.
